% 국문 이력서 조판 선언 — resume-ko.md 전용.
% cover-letter의 여유 있는 행간 대신 훑는 문서에 맞춘 밀도를 둔다.
\usepackage{xetexko}
\setmainhangulfont{Pretendard}[Ligatures=TeX]
\setsanshangulfont{Pretendard}
\setmonohangulfont{D2Coding}
\setmainhanjafont{D2Coding}

\usepackage[dvipsnames]{xcolor}
\definecolor{accent}{HTML}{1B3A57}
\definecolor{muted}{HTML}{5A6672}
\definecolor{rulegray}{HTML}{B9C2CC}
\usepackage[colorlinks=true,allcolors=accent,breaklinks=true,pdfborder={0 0 0}]{hyperref}

\usepackage{titlesec}
\titleformat{\section}{\normalfont\large\bfseries\color{accent}}{}{0pt}{}[\vspace{1pt}\color{rulegray}\titlerule]
\titlespacing*{\section}{0pt}{8pt plus 2pt minus 1pt}{3pt}
\titleformat{\subsection}{\normalfont\normalsize\bfseries}{}{0pt}{}
\titlespacing*{\subsection}{0pt}{6pt plus 1pt minus 1pt}{2pt}
\titleformat{\subsubsection}{\normalfont\small\bfseries\color{muted}}{}{0pt}{}
\titlespacing*{\subsubsection}{0pt}{4pt plus 1pt minus 1pt}{1pt}

\usepackage{enumitem}
\setlist[itemize]{leftmargin=1.15em,itemsep=1pt,topsep=1.5pt,parsep=0pt,partopsep=0pt}
\setlength{\parindent}{0pt}
\setlength{\parskip}{2.2pt plus 1pt minus 0.4pt}
\linespread{1.03}
\setlength{\emergencystretch}{1.2em}

\usepackage{fancyhdr}
\pagestyle{fancy}
\fancyhf{}
\renewcommand{\headrulewidth}{0pt}
\fancyfoot[L]{\footnotesize\color{muted}김정한 (Junghan Kim) · [email removed]}
\fancyfoot[R]{\footnotesize\color{muted}\thepage}
